\section{\textit{kunnen} (poder, capacidad)}
\textbf{Significado.} Expresa habilidad o posibilidad realista.

\begin{tabular}{@{}lll@{}}
\toprule
Persona & Forma & Nota \\
\midrule
ik & \textit{kan} &  \\
jij/je & \textit{kan} o \textit{kunt} & ambas válidas en A1 \\
wij/jullie/zij & \textit{kunnen} & plural regular \\
\bottomrule
\end{tabular}

\textbf{Ejemplos}
\begin{itemize}[leftmargin=*]
  \item \textit{Ik kan zwemmen.} (Sé nadar.)
  \item \textit{Wij kunnen niet komen.} (No podemos venir.)
  \item Pregunta: \textit{Kan je morgen om acht uur?}
\end{itemize}

\textbf{Mini práctica}
\begin{itemize}[leftmargin=*]
  \item Completa con la forma correcta: \textit{Hij \underline{\hspace{1.2cm}}\underline{\hspace{1.2cm}}\underline{\hspace{1.2cm}} (kunnen) Nederlands praten}. 
  \item Haz una pregunta de disponibilidad con \textit{kan}. 
\end{itemize}
